% Automatically generated by R/30_make_posthoc_table_and_figure.R
% Generated as a standard tabular environment; no tabularx dependency.
\begin{table}[t]
\centering
\footnotesize
\caption{Post-hoc examples used to illustrate the information tiers. The table records the source, support setting, definition of reported or reconstructed $R^2$, and information available for adjustment; quantitative changes in the metrics are displayed in Fig.~\ref{fig:posthoc-examples-summary}. Empty or unavailable quantities in the underlying results indicate that the corresponding tier is not identifiable from the published information.}
\label{tab:posthoc-examples-overview}
\begin{tabular}{p{0.14\textwidth} p{0.10\textwidth} p{0.18\textwidth} p{0.29\textwidth} p{0.07\textwidth} p{0.15\textwidth}}
\hline
Source & Variable & Prediction/reference support & Information used & Tier & Main limitation \\
\hline
\citet{zandler2015dwarf} & Total biomass & 30 m map pixel / 60 m field site; two 4 m x 4 m subplots & Spatial-CV RMSE; approximate response SD reconstructed from RMSE/SD ratio; allometric and within-site reference-error components; $q_{\varepsilon}$ = 0.172, $\lambda$ = --; predictive\_skill & Tier 1 & Reference support is a field-site estimate from sparse subplots; response variance is reconstructed from a rounded RMSE/SD ratio; no semivariogram is available. \\
\citet{brenning2005bts} & BTS & 30 m grid cell / BTS point observation & ordinary-regression R2, residual semivariogram, transferred measurement-error variance; $q_{\varepsilon}$ = 0.012, $\lambda$ = 0.806; squared\_correlation & Tier 2 & Measurement-error variance is not Zermatt-specific; it is transferred from related Norway BTS data reported in the source study. \\
Meuse data$^{a}$ & Log cadmium & 100 m grid cell / treated as point soil sample & LOOCV predictive R2, response variance and fitted response semivariogram; $q_{\varepsilon}$ = 0.000, $\lambda$ = 0.575; predictive\_skill & Tier 3 & Bulk soil samples treated as point support for illustration. \\
Meuse data$^{a}$ & Log zinc & 100 m grid cell / treated as point soil sample & LOOCV predictive R2, response variance and fitted response semivariogram; $q_{\varepsilon}$ = 0.000, $\lambda$ = 0.806; predictive\_skill & Tier 3 & Bulk soil samples treated as point support for illustration. \\
NO$_2$ Germany 2018 & Annual mean NO2 & 2 km grid cell / monitoring station annual mean & station data; simple 10-fold CV metrics; fitted response semivariogram; sigma\_epsilon\textasciicircum{}2 set to 0; $q_{\varepsilon}$ = 0.000, $\lambda$ = 0.394; squared\_correlation & Tier 3 & Illustrative validation model; annual-mean measurement error set to zero; fitted nugget not decomposed. \\
\citet{vizcaino2018no2} & Annual mean NO2 & 100 m grid cell / AirBase monitoring station & station data; recovered AirBase/EEA coordinates; response variance; CV RMSE; fitted response semivariogram; nugget partition scenario; $q_{\varepsilon}$ = 0.278, $\lambda$ = 0.445; predictive\_skill & Tier 2 sens. & Does not reproduce pan-European raster map; fitted nugget decomposition is a sensitivity assumption. \\
\citet{vizcaino2018no2} & Annual mean NO2 & 100 m grid cell / AirBase monitoring station & station data; recovered AirBase/EEA coordinates; response variance; CV RMSE; fitted response semivariogram; nugget partition scenario; $q_{\varepsilon}$ = 0.000, $\lambda$ = 0.445; predictive\_skill & Tier 2 sens. & Does not reproduce pan-European raster map; fitted nugget decomposition is a sensitivity assumption. \\
\citet{balenzano2021sentinel1} & Soil moisture & Sentinel-1 soil moisture product, approximately 1 km / in situ point-scale soil moisture at 0.05 m depth & reported RMSE, ubRMSE, bias, SRE, intrinsic RMSE, and OLS/WLS correlations; $q_{\varepsilon}$ = --, $\lambda$ = --; not\_used\_reported\_sre\_example & Reported SRE & Direct SRE example; not a covariance-model post-hoc adjustment. \\
\citet{zhao2021subsurface} & Soil moisture & SMAP 9 km pixel / CRNS footprint & r, ubRMSD, TC-estimated CRNS error SD; $q_{\varepsilon}$ = --, $\lambda$ = --; NA & Partial Tier 1 & Reference variance not tabulated; only ubRMSD can be corrected. \\
\citet{plattner2006swe} & SWE & 10 m grid cell / local snow survey site & OLS R2, response SD, range, nugget-like term, two measurement-error decompositions; $q_{\varepsilon}$ = 0.000--0.015, $\lambda$ = 0.522--0.885; squared\_correlation & Tier 2 sens. & OLS R2 used as squared-correlation proxy; nugget split into measurement error and true microscale variability is heuristic. \\
\hline
\multicolumn{6}{p{0.96\textwidth}}{\footnotesize $^{a}$For predictive/test-set $R^2$, adjustments are performed on the MSE scale and then transformed back to $R^2$ using the target-support variance. For the SWE sensitivity rows, the ordinary-regression $R^2$ is treated as a squared-correlation proxy; the nugget-like term is split heuristically into measurement error and true microscale variability. The S1-SM reported-SRE row uses the spatial representativeness error adjustment as reported in the original Sentinel-1 soil-moisture validation study. The SMAP--CRNS row uses the TC-estimated CRNS reference-error standard deviation for a partial Tier~1 ubRMSD adjustment only. For the NO2--EU rows, station coordinates recovered from AirBase metadata are used to fit a European semivariogram, with sensitivity rows that decompose the fitted nugget into measurement error and true microscale variation. The Meuse log-zinc and log-cadmium examples use LOOCV predictive $R^2$ and treat bulk soil-sample observations as point-support observations for illustration, although they have finite sampling support. If the inferred support-error variance exceeds the observed prediction-error variance, the Tier~2 RMSE adjustment is flagged as non-informative rather than truncated to zero.} \\
\end{tabular}
\end{table}
